% Generated draft: requires mathematical and source review.
\documentclass[11pt]{article}
\usepackage[margin=25mm]{geometry}
\usepackage{fontspec}
\setmainfont{Latin Modern Roman}
\usepackage{amsmath,amssymb,mathtools,unicode-math}
\setmathfont{Latin Modern Math}
\usepackage{enumitem,needspace,xurl,hyperref}
\setlength{\parindent}{0pt}
\setlength{\parskip}{5pt}
\setlength{\emergencystretch}{3em}
\newcommand{\N}{\mathbb{N}}
\newcommand{\Z}{\mathbb{Z}}
\newcommand{\Q}{\mathbb{Q}}
\newcommand{\R}{\mathbb{R}}
\newcommand{\C}{\mathbb{C}}
\newcommand{\F}{\mathbb{F}}
\newcommand{\A}{\mathbb{A}}
\newcommand{\PP}{\mathbb{P}}
\newcommand{\E}{\mathbb{E}}
\newcommand{\eps}{\varepsilon}
\newcommand{\dd}{\,\mathrm{d}}
\begin{document}
\section*{UnsolvedMath ID: 801018}
\textbf{Source problem code:} catalogue entry
\subsection{Definitions and mathematical statement}
Write \(\omega=\{0,1,2,\ldots\}\). For a tree \(T\), let \([T]\) denote its set of infinite paths. The supplied record does not specify the ambient class of trees or define the adjective ``positive.'' Accordingly, retain \(\operatorname{Pos}(T)\) as an unresolved predicate expressing the source's intended positivity condition. With these qualifications, the supplied principle is
\[
P^-:\quad \forall T\,\Bigl(\operatorname{Pos}(T)\ \longrightarrow\ \exists (X_n)_{n\in\omega}\,\bigl[(\forall n\in\omega\;X_n\in[T])\ \land\ (\forall m,n\in\omega\;(m\ne n\longrightarrow X_m\ne X_n))\bigr]\Bigr).
\]
Here the existential quantifier asks for a single sequence of pairwise distinct paths. Its precise coding in second-order arithmetic depends on the missing source definitions.

The target question is whether
\[
\mathrm{RCA}_0\vdash (P^-\longrightarrow P),
\]
where \(\mathrm{RCA}_0\) is the usual recursive comprehension base theory of second-order arithmetic. Crucially, the supplied excerpt does not define \(P\); it also uses an undefined principle \(P^+\). Thus this is an explicitly incomplete formulation of the intended provability question, and no definition of either principle is inferred here.

For context, the supplied excerpt reports the following relationships:
\[
\mathrm{WKL}_0\vdash P^+,
\]
\[
\mathrm{RCA}_0\vdash P^+\to P\to P^-\to\mathrm{1\text{-}RAN}\leftrightarrow\mathrm{WWKL}_0,
\]
and
\[
P^+\nvdash\mathrm{WKL}_0,\qquad
\mathrm{RCA}_0\nvdash(P\to P^+),\qquad
\mathrm{RCA}_0\nvdash(\mathrm{1\text{-}RAN}\to P^-).
\]
Here \(\mathrm{WKL}_0\) and \(\mathrm{WWKL}_0\) denote the usual weak and weak weak K\"onig's lemma systems. These displays preserve the source's notation and are supplied background claims, not independently verified assertions. The excerpt glosses \(\mathrm{1\text{-}RAN}\) as ``there exists a 1-random,'' but does not specify its relativization or formal quantifiers. It also does not specify the background theory implicit in \(P^+\nvdash\mathrm{WKL}_0\). These omissions must be resolved before the background relationships or the target can receive a fully precise formal interpretation.
\subsection{Short English statement}
The supplied question asks whether \(\mathrm{RCA}_0\) proves that every positive tree having a sequence of pairwise distinct infinite paths, as expressed by \(P^-\), implies the principle \(P\). The excerpt does not define \(P\) or positivity, so the intended implication cannot yet be stated completely. The registry labels the problem open; this is a supplied status, not an independently verified assessment.
\subsection{Sources}
\begin{itemize}
\item Local UnsolvedMath registry: supplied complete record with permanent ID \(801018\), exact source problem code \texttt{catalogue entry}, title ``Open Problems in Mathematical Logic --- Problem 30,'' in the AMR Open Problem Lists collection, Logic category. The registry attributes the record to Barmpalias, Gao et al., gives proposed year 2026, and labels its status \texttt{open}.
\item Supplied source attribution: Barmpalias, Gao et al., \emph{Open Problems in Mathematical Logic} (2026), Problem 30, \url{https://arxiv.org/abs/2608.26628}. The source metadata reports access on 2026-10-16 and extraction by \texttt{arxiv\_latex}. These are supplied provenance details; the source and its mathematical claims have not been independently checked here.
\end{itemize}
\end{document}
